\section{问题重述与建模边界}

本演示研究一个小规模、确定性的资源分配问题。设候选物品集合含 $n$ 个物品，记为 $I$。
物品 $i$ 的成本为 $c_i>0$，价值为 $v_i\geq 0$，总容量为
$C>0$。需要决定是否选择每个物品，使所选物品成本不超过容量且总价值尽可能高。
该结构是经典 $0$--$1$ 背包形式，但本文只解释 D4 合成运行包中的有限实例，不把它外推为真实业务模型。

发布模块不重新选择算法，也不改变实验口径。本文所用路线、正式运行标识
\texttt{\ExperimentRunId}、样本状态和 Metric 均来自冻结引用；运行包本身作为开发资料引用
\cite{D4SyntheticRun}。第三个实例的失败是有意注入的测试情形，目的在于验证部分结果不会被排版过程隐藏。

\subsection{基本假设}

为使算法和指标闭合，采用以下有限假设：
\begin{enumerate}
  \item 每个物品至多选择一次，选择变量只取 $0$ 或 $1$；
  \item 同一实例内成本和容量使用一致的抽象单位，价值使用 \texttt{value}；
  \item 物品价值可加，且不考虑物品之间的相互作用；
  \item 只对状态为 \texttt{completed} 的完整观测计算均值，失败实例保留在覆盖率分母中；
  \item 穷举结果仅在给定有限物品集合的全部子集内称为最优，不代表大规模问题的可行求解策略。
\end{enumerate}

\subsection{符号说明}

\begin{table}[H]
  \centering
  \caption{主要符号及其含义}
  \begin{tabularx}{\linewidth}{cYc}
    \toprule
    符号 & 含义 & 单位\\
    \midrule
    $I$ & 当前实例的候选物品集合 & --\\
    $c_i$ & 物品 $i$ 的成本 & cost\\
    $v_i$ & 物品 $i$ 的价值 & value\\
    $C$ & 当前实例容量上限 & cost\\
    $x_i$ & 是否选择物品 $i$ 的二元变量 & --\\
    $V_G,V_E$ & 密度贪心与穷举法得到的目标值 & value\\
    $r_i=v_i/c_i$ & 物品 $i$ 的单位成本价值密度 & value/cost\\
    \bottomrule
  \end{tabularx}
\end{table}
